\chapter{Learning, structural adaptation and evolution}
A model can change its answer without changing its weights. It can change
its weights without changing its architecture. It can change its architecture
without becoming a population undergoing selection. Calling all three
processes \emph{adaptation} is acceptable in conversation, but inadequate
for implementing an experiment.

We will build a small structural search whose task is deliberately
transparent. Two observed features agree during training. Only one remains
predictive when that agreement is broken. The learner fits the training
set almost perfectly using either feature. Can a search over feature
masks identify the robust representation? Which data choose the winner,
and which pieces of a parent's implementation does a child inherit?

Prerequisites are gradients and evaluation from Chapters 2--4, transient
state from Chapter 6, and generated graphs from Chapter 12. The example is
a deterministic two-feature regression problem, not evidence of genomic
capability or a general architecture-search advantage.

\section{Four objects and four update contracts}
\index{state update}\index{structural proposal}
Write a candidate as $(G,\theta)$, where $G$ is a computational structure
and $\theta$ its trainable parameters. Let $h_t$ denote transient state
during an input sequence and $\mathcal P_g$ a population at generation $g$.
A useful separation is
\begin{align}
 h_{t+1}&=f_{G,\theta}(h_t,x_t), &&\text{state update},\\
 \theta_{k+1}&=U_G(\theta_k,D_{\mathrm{train}}),&&\text{parameter learning},\\
 G'&=M(G,\xi),&&\text{structural proposal},\\
 \mathcal P_{g+1}&=\mathcal S(\mathcal P_g,\mathcal C_g,E),&&\text{selection}.
 \label{eq:clocks}
\end{align}
Here $t$ is an input step, $k$ an optimization step, $\xi$ a proposal choice,
$\mathcal C_g$ candidate offspring and $E$ an evaluation procedure.
These indices name causal boundaries; they need not correspond to
universally separated wall-clock times.

\Plate{01-clocks}{State, parameters, structure and population have distinct
owners and update events. A state trajectory occurs under one parameter
version; a training run occurs under one structural contract. The example
uses explicit resets at these boundaries.}{fig:clocks}

An attention cache changing during inference is not parameter learning.
Turning an expert on for one token is not necessarily structural learning:
the graph may already contain all experts and a fixed routing rule.
A rewritten graph is a structural proposal even before anyone has measured
its usefulness. Evolution requires an inheritance/variation/selection
process over candidates, not merely an operation named \texttt{mutate}.

Biological inheritance motivates vocabulary but does not license an
identity claim. A software mask has no cell machinery, metabolism or
reproductive ecology. We will use \Term{genotype} only as a declared
encoding of a computational candidate, with its interpretation specified.

\section{A two-feature problem that exposes the difference}
\index{gradient descent}\index{spurious feature}
Take training rows $(-1,-1)$ and $(1,1)$, with target $y=2x_1$.
The validation set contains all four sign combinations
$(x_1,x_2)\in\{-1,1\}^2$, still with $y=2x_1$.
All features and targets are dimensionless synthetic values.

A binary mask $m=(m_1,m_2)$ defines the active structure:
\begin{equation}
 \widehat y=x^\top(m\odot w),\qquad
 L_{\mathrm{train}}(w;m)=\frac1N\|X(m\odot w)-y\|_2^2.
 \label{eq:model}
\end{equation}
$X\in\mathbb R^{N\times2}$, $w,m\in\mathbb R^2$ and $y\in\mathbb R^N$.
The mask is held fixed during gradient descent. We do not differentiate
through its discrete selection.

\Plate{02-gradient}{Training samples lie on $x_2=x_1$, so either coordinate
can explain them. The validation grid breaks that correlation while
retaining $y=2x_1$. The experiment tests a specified shift, not all possible
future distributions.}{fig:features}

Let $r=X(m\odot w)-y$. Since changing $w_j$ changes predictions by
$m_jX_{:,j}$, differentiation gives
\begin{equation}
 \nabla_wL=\frac2N m\odot(X^\top r).
 \label{eq:gradient}
\end{equation}
An inactive coordinate receives zero gradient. Starting all weights at
zero therefore keeps inactive entries zero. That storage convention avoids
hidden inherited values becoming active after a later mask change.

The reference uses 40 full-batch updates with learning rate $\eta=0.1$.
For a fixed mask the Hessian is
$H=(2/N)\operatorname{diag}(m)X^\top X\operatorname{diag}(m)$.
On every positive-eigenvalue mode, gradient descent contracts when
$0<\eta<2/\lambda_{\max}(H)$. Zero-eigenvalue directions are unidentifiable
and retain their initialization. For mask 11, $\lambda_{\max}=4$;
the chosen rate satisfies the bound.

\Listing{fit}

\section{Derive a numerical oracle without running the loop}
For a single active feature, the training loss is $(w-2)^2$.
One update is $w_{k+1}=w_k-0.2(w_k-2)=0.8w_k+0.4$.
With $w_0=0$, solve the geometric recurrence:
\begin{equation}
 w_k=2(1-0.8^k).
 \label{eq:oneweight}
\end{equation}
Either feature obtains the same training loss because the training columns
are identical.

For mask 11, symmetry and zero initialization keep $w_1=w_2=a$.
Each component's gradient is $2(2a-2)$, giving
$a_{k+1}=0.6a_k+0.4$ and
\begin{equation}
 a_k=1-0.6^k.
 \label{eq:twoweights}
\end{equation}
These formulas are independent oracles for the iterative implementation,
including its gradient scaling and initialization.

On the validation grid, $\mathbb E[x_1x_2]=0$ and
$\mathbb E[x_1^2]=\mathbb E[x_2^2]=1$. Hence
\begin{equation}
 L_{\mathrm{valid}}=(m_1w_1-2)^2+(m_2w_2)^2.
 \label{eq:validation}
\end{equation}
At convergence, masks 10, 01 and 11 have validation losses 0, 8 and 2.
The all-zero predictor has loss 4. Perfect training fit does not identify
the causal feature; the deliberately chosen validation intervention supplies
information absent from the training set.

We select by
\begin{equation}
 J(m)=L_{\mathrm{valid}}(w_{40}(m);m)+0.05\|m\|_0.
 \label{eq:score}
\end{equation}
The second term penalizes active features, not weight magnitude. Its units
match squared target units per active feature. It is a declared preference,
not a calibrated hardware cost.

\begin{center}\input{\Generated/candidate-table.tex}\end{center}

The dense representation fits the training set better at step 40, yet the
one-feature 10 mask wins the declared validation objective. Nothing here
establishes that sparse models generally outperform dense ones.

\section{Selection needs a reproducible lineage}
\index{search!budget}
Our \Term{elitist search} retains the best candidate among a parent and
its one-bit mutations. Generation zero evaluates only parent 01.
At the next generation, candidates are 01, 00 and 11; 11 wins.
Its next neighborhood contains 11, 01 and 10; 10 wins.
All weights are freshly trained from zero for every evaluation, including
the parent. Ties use lexicographic mask order, never completion order.

\Listing{select}
\Plate{03-lineage}{The lineage inherits masks only. Each candidate receives
the same 40-step training recipe, followed by the same validation objective.
Dashed descent edges encode structural inheritance; no trained weight
crosses them in this experiment.}{fig:lineage}

\begin{center}\input{\Generated/history-table.tex}\end{center}
\Listing{evolve}

This is a $(1+\lambda)$ population-selection scheme with one survivor,
not a diverse many-survivor population. The parent remains eligible,
so deterministic reevaluation under a fixed objective cannot worsen the
best score. Noise, changing data or changed evaluation budgets break that
argument. A monotone validation score also says nothing about test
performance.

Equal step counts are not equal compute. A candidate with two active
features performs more useful scalar products than one with one feature;
our NumPy representation still stores both columns. This reference isolates
optimization budget and semantics, not actual sparse-kernel efficiency.
A serious comparison should report candidate evaluations, training steps,
data exposures, active parameter count, measured compute and wall time.
Caching repeated candidates could save work, but the cache key must include
data, initialization, training recipe and code version, not just the mask.

\section{Inheritance is a mapping, not a file copy}
\index{inheritance!parameter map}
An alternative policy can inherit trained weights. For a stable feature
coordinate $i$, define
\begin{equation}
 w_i'=\begin{cases}
 w_i,&m_i=m_i'=1,\\
 0,&\text{otherwise}.
 \end{cases}
 \label{eq:inherit}
\end{equation}
A newly enabled feature starts at zero; a disabled feature's stored weight
is cleared. The policy needs a stable meaning for coordinate $i$.
Equal vector lengths are not enough if features were reordered or replaced.

\Listing{inherit}
\Plate{05-inheritance}{The inherited coordinate keeps its value, the newly
enabled coordinate starts at zero, and transient sequence state resets.
This is an alternative policy, not the policy used to generate the search
table. Optimizer state would require a separate inheritance rule.}{fig:inheritance}

Weight inheritance changes the experiment: candidates no longer begin
from an equally trained state. It may accelerate search, but inherited
quality and structural quality become coupled. Momentum buffers, adaptive
optimizer statistics, random-number state, normalization statistics and
data order must also be addressed. Copying a weight without its optimizer
history is a different algorithm from copying both.

In terminology often used for computational evolution, a \Term{Baldwinian}
evaluation can let within-candidate learning affect fitness without
inheriting its learned weights; a \Term{Lamarckian} variant can transmit them.
These labels are useful policy shorthand here, not claims about biological
mechanisms. Always publish the actual state map.

\section{Search uses up evaluation information}
\index{validation!adaptive reuse}\index{multiple testing}
The validation set participates in selection and is therefore part of
the search process. It is not an untouched final test. Our complete
four-point grid is an exhaustive check of one synthetic condition;
we report no separate generalization estimate.

For intuition, suppose $K$ independent null candidates each have probability
$\alpha$ of a misleading apparent improvement. Then
\begin{equation}
 \Pr(\text{at least one false positive})=1-(1-\alpha)^K.
 \label{eq:multiple}
\end{equation}
At $\alpha=0.05$ and $K=20$, this is approximately 0.642.
The calculation is a model of repeated testing, not a calibrated error
rate for adaptive architecture search. Candidates and their errors are
usually dependent; repeated querying can also make later candidates
depend on earlier validation results.

\Plate{04-budget}{A larger search increases the chance of finding an
apparently favorable null result under the stated independent-test model.
Training, selection and final assessment must retain distinct roles.
The curve is illustrative probability arithmetic, not measured search
performance.}{fig:searchbudget}

Use an untouched final set only after freezing the selection procedure.
For genomic tasks, split by the scientifically relevant unit rather than
assuming random windows are independent; related sequences and shared
sources can leak information. Predeclare the task, grouping rule, baselines,
search budget, stopping rule and uncertainty analysis. Repeating seeds
does not repair a systematically leaked split.

\section{A current frontier: evaluator-driven program evolution}
AlphaEvolve's 2025 technical report describes an LLM proposing program
changes, automated evaluation, and an evolutionary database that balances
reuse and diversity \cite{alphaevolve}. Its optional evaluation cascade
rejects unpromising programs before more costly stages. This is a relevant
systems pattern, not the algorithm implemented by our two-bit reference.

The important engineering lesson is that the evaluator is part of the
search system's specification. A candidate can exploit an incomplete metric,
a timeout policy, a data leak or an implementation bug. Passing tests is
evidence about those tests, not an unrestricted correctness proof.
For executable candidate code, isolation, resource limits and evaluator
integrity are separate requirements; this chapter's safe numeric masks
avoid executing generated source.

A useful next experiment would compare a fixed model, random mask search
and neighborhood evolution under equal total evaluation budgets.
Report both best validation score and final held-out outcome. Include failed
candidates and search overhead. Without that comparison, an attractive
lineage diagram is not evidence that evolution was necessary.

\section{Exercises with worked answers}
\begin{enumerate}
\item \textbf{Classify.} A recurrent state changes while weights are frozen.
\textit{Answer:} This is a state update, not parameter or structural learning.
\item \textbf{Derive.} Differentiate Equation~\ref{eq:model}.
\textit{Answer:} Apply the chain rule to each residual:
$\partial L/\partial w_j=(2/N)\sum_nr_nX_{nj}m_j$.
Stacking coordinates gives Equation~\ref{eq:gradient}.
\item \textbf{Oracle.} Compute the single-feature weight after one and two
steps. \textit{Answer:} $w_1=0.4$ and $w_2=0.72$,
matching $2(1-0.8^k)$.
\item \textbf{Identifiability.} Why is mask 11's optimum not unique on the
training set? \textit{Answer:} Only $w_1+w_2$ matters there. Every pair
summing to 2 minimizes training loss; symmetric initialization and updates
select the symmetric solution.
\item \textbf{Shift.} Derive the spurious feature's limiting validation loss.
\textit{Answer:} Predictions are $2x_2$; the mean squared error is
$4\mathbb E[(x_2-x_1)^2]=8$ on the independent sign grid.
\item \textbf{Budget.} Does forty steps per mask imply equal FLOPs?
\textit{Answer:} No. Active work, tensor representation and execution kernels
can differ. The controlled variable is the number of updates.
\item \textbf{Lineage.} Why does the first mutation round not choose 10?
\textit{Answer:} It differs from parent 01 in two bits and is not in the
one-bit neighborhood. The search reaches it through 11.
\item \textbf{Inheritance.} Map $(1,0)$ with stored weights $(2,9)$ to $(1,1)$.
\textit{Answer:} Equation~\ref{eq:inherit} yields $(2,0)$;
the stale inactive value 9 must not silently become an inherited skill.
\item \textbf{Selection proof.} When is best score monotone?
\textit{Answer:} When the parent is included and all candidates are scored
deterministically under the same objective. Changed conditions invalidate
the comparison; noisy minima can improve without true quality improving.
\item \textbf{Multiple testing.} For $K=2$ and $\alpha=.05$, compute the
familywise probability. \textit{Answer:} $1-.95^2=.0975$,
not .05. The independence assumption must be stated.
\item \textbf{Reproducibility.} Specify a candidate cache key.
\textit{Rubric:} Include structural encoding, data/split identity,
initialization and seed, optimization budget, hyperparameters, evaluator
version, numeric precision and code/environment version.
\item \textbf{Research design.} Test whether evolution is useful.
\textit{Rubric:} Compare fixed, random-search and evolutionary baselines;
match total budgets; freeze selection before final testing; report multiple
runs, failed candidates, overhead and uncertainty. Do not tune on final data.
\end{enumerate}

\section{From adaptation policies to typed candidates}
\textbf{Reproduce the chapter.} From the repository root, run
\texttt{python3 drafts/ch13/build.py --assets-only}, then
\texttt{python3 -m pytest tests/test\_ch13\_reference.py}.
The first command regenerates \texttt{results.json}, tables, plots and literal
code listings; the tests compare the implementation with independent
oracles and explicit failure cases. Run the builder without
\texttt{--assets-only} to compile this review PDF.

We can now distinguish state transition, learning, structural proposal and
selection by inspecting which object changes and which evidence authorizes
the change. The synthetic example demonstrates an evaluation failure mode,
not a universal recipe for robust models.

\textbf{Working vocabulary.} A genotype encodes a candidate; phenotype is
its interpreted computation; inheritance maps parent state to child state;
fitness is an evaluator-defined quantity; elitism retains a qualifying
parent; a holdout ceases to be untouched once it guides decisions.
Chapter 14 makes candidate structure explicit enough to validate before
execution: typed computational genes, owned state and an expression trace.
